\label{conclusion}

\textbf{A. Tổng kết đóng góp khoa học}

Đề tài đã hoàn thành việc xây dựng hệ thống nhận dạng hoạt động lính cứu hỏa theo thời gian thực (RFAR) dựa trên nền tảng TinyML, bao gồm bốn đóng góp khoa học chính.

Thứ nhất, một tập dữ liệu chuyên biệt mô phỏng 11 hoạt động đặc thù của lính cứu hỏa đã được xây dựng từ 20 tình nguyện viên trong điều kiện diễn tập thực tế, với tổng thời gian thu thập 371 phút. Tập dữ liệu này là đóng góp có giá trị cho cộng đồng nghiên cứu HAR trong bối cảnh thiếu hụt dữ liệu công khai cho lính cứu hỏa.

Thứ hai, quy trình xử lý tín hiệu mới kết hợp kỹ thuật cửa sổ trượt 3 giây với phát hiện chỉ số bất thường AIx đã được đề xuất và kiểm chứng. Phương pháp này hiệu quả hơn cửa sổ trượt cố định thông thường trong việc xử lý các chuyển tiếp hoạt động và phát hiện sự kiện ngã, giảm thiểu phân loại nhầm do chồng lấp hoạt động.

Thứ ba, hệ thống RFAR với mô hình RF tối ưu (17 cây, độ sâu 14, kích thước 0,92~MB) đã được triển khai thành công trên vi điều khiển ESP32 với độ trễ xấp xỉ 9,6~ms. Hệ thống đạt F1$_{mi} = 96{,}2\%$ trong thực nghiệm thời gian thực và F1$_{mi} \geq 96{,}6\%$ trên hai tập dữ liệu công khai, khẳng định tính khả thi của giải pháp TinyML cho ứng dụng HAR trong môi trường khẩn cấp.

Thứ tư, thiết bị đeo được thiết kế và chế tạo với trọng lượng dưới 200~g, khả năng hoạt động liên tục 24~giờ và phạm vi truyền thông vô tuyến 8000~m, đáp ứng yêu cầu thực tế trong các tình huống cứu hộ.

\textbf{B. Đánh giá tiến độ thực hiện đề tài}

Đề tài đã hoàn thành các công việc cốt lõi của giai đoạn nghiên cứu hiện tại. Cụ thể, các công việc đã hoàn thành bao gồm: (1) thu thập và xử lý tập dữ liệu lính cứu hỏa từ 20 tình nguyện viên; (2) thiết kế và chế tạo thiết bị đeo nhỏ gọn tích hợp ADXL345 và ESP32 hoạt động trên 24~giờ; (3) phát triển quy trình xử lý tín hiệu mới với cơ chế AIx; (4) tối ưu hóa và triển khai mô hình RF trên vi điều khiển; (5) đánh giá hệ thống trên bốn tập dữ liệu và xác nhận hiệu suất trong thực nghiệm thực tế.

Kết quả đạt được phản ánh tiến độ tích cực và vượt kỳ vọng ban đầu ở khía cạnh chất lượng thực nghiệm và độ trễ suy luận trên thiết bị nhúng. Bài báo mô tả toàn bộ hệ thống RFAR đã được chấp nhận đăng trên tạp chí IEEE Sensors (JCR Q1) vào tháng 7 năm 2025, thể hiện chất lượng nghiên cứu được ghi nhận tại diễn đàn khoa học quốc tế.

\textbf{C. Kế hoạch công việc tiếp theo}

Trong các giai đoạn tiếp theo, nghiên cứu sẽ tập trung vào ba hướng chính: (1) tối ưu hóa thiết kế phần cứng, hướng đến tích hợp thiết bị vào trang phục bảo hộ tiêu chuẩn của lính cứu hỏa nhằm nâng cao tính thực tế trong triển khai thực địa; (2) mở rộng đối tượng thu thập dữ liệu với sự tham gia của lính cứu hỏa nữ và người có đặc điểm thể chất đa dạng hơn; và (3) tích hợp hệ thống nhận dạng hoạt động với module định vị trong nhà và cơ chế phát cảnh báo khẩn cấp tự động để hoàn thiện giải pháp hỗ trợ toàn diện cho đội cứu hộ.

\textbf{D. Kiến nghị và đề xuất}

Để tạo điều kiện thuận lợi cho việc triển khai hệ thống RFAR trong bối cảnh thực tế, nhóm nghiên cứu đề xuất: (1) hỗ trợ kinh phí và điều kiện thực nghiệm mở rộng tại nhiều cơ sở phòng cháy chữa cháy trong cả nước; (2) tạo điều kiện hợp tác chính thức với lực lượng cảnh sát phòng cháy chữa cháy để thu thập dữ liệu trong các buổi diễn tập quy mô lớn; và (3) kết nối đề tài với chương trình nghiên cứu phát triển hệ thống định vị trong nhà tổng thể dành cho lính cứu hỏa đang được triển khai tại Trường Đại học Phenikaa. Những điều kiện này sẽ tạo nền tảng để chuyển hóa kết quả nghiên cứu thành sản phẩm ứng dụng thực tế có giá trị xã hội cao.